\documentclass[conference]{IEEEtran}
\usepackage{booktabs}
\usepackage{amsmath}
\usepackage{graphicx}
\usepackage{multirow}
\usepackage{url}

\title{Bias Concentration Is Not Bias Expression:\\A Prefix-Dimension Study of Matryoshka Representation Learning Embeddings}

\author{\IEEEauthorblockN{Anonymous Authors}}

\begin{document}
\maketitle

\begin{abstract}
Matryoshka Representation Learning (MRL) enables dynamic prefix truncation of embedding vectors, but its social-bias behavior across prefix dimensions remains underexplored. We study gender bias as the primary axis (with race as exploratory evidence) across MPNet-Standard, MPNet-MRL, Nomic-MRL, mxbai-MRL, Qwen3-0.6B-MRL, and Qwen3-4B-MRL. We evaluate standard bias metrics across dimensions (WEAT, SEAT, StereoSet, CrowS-Pairs) and geometric concentration $C(d)$, then add a variance-based view $C_{var}(d)$ with bootstrap confidence intervals. We further define the anisotropy factor $\alpha(d)=C_{var}(d)/\|v_{1:d}\|^2$ to separate geometric bias mass from statistically expressed bias variance. The key finding is that early dimensions can appear bias-concentrated geometrically while simultaneously underexpressing bias variance ($\alpha(d)<1$), implying that concentration and expression are not equivalent. This provides a clearer measurement framework for fairness analysis in dimension-adaptive retrieval systems.
\end{abstract}

\begin{IEEEkeywords}
matryoshka representation learning, embedding bias, fairness, anisotropy, dimensional truncation, social bias
\end{IEEEkeywords}

\section{Introduction}
MRL embeddings are designed for prefix truncation while preserving task utility, making them attractive for latency-sensitive retrieval. However, social bias under truncation is insufficiently characterized. Existing bias analyses usually inspect full-dimension embeddings or metric trends without separating geometric concentration from variance expression. This work targets that gap.

\textbf{Central claim.} Bias concentration in early dimensions is not the same as bias expression in early dimensions. In MRL models, geometric bias mass may be front-loaded while expressed variance is suppressed.

\textbf{Scope.} Gender bias is the primary analysis axis. Race is included only as secondary/exploratory evidence from earlier-stage metrics.

\section{Research Questions}
\textbf{RQ1:} How do standard bias metrics change as prefix dimension shrinks in MRL embeddings?\\
\textbf{RQ2:} Is bias geometrically concentrated in early dimensions?\\
\textbf{RQ3:} Is geometric concentration actually expressed as statistical variance, or suppressed by anisotropy/covariance structure?

\section{Related Work}
This study connects prior work on: (i) word/sentence embedding bias tests (WEAT, SEAT), (ii) benchmark-based stereotype preference (StereoSet, CrowS-Pairs), and (iii) anisotropy and representation geometry in high-dimensional embedding spaces. It also extends MRL evaluation from utility-centric analysis to fairness-aware dimension-wise analysis.

\section{Problem Setup and Metrics}
Let $f^{(d)}(x)$ denote the first $d$ dimensions of embedding $f(x)$, where $D$ is the full embedding dimension and $d\leq D$. Let $v$ be a full-dimension, unit-normalized bias direction ($\|v\|=1$) estimated from attribute groups (e.g., male vs female).

\subsection{Standard Metrics (RQ1)}
WEAT, SEAT, StereoSet, and CrowS-Pairs are computed across tested dimensions per model.

\subsection{Geometric Concentration (RQ2)}
\begin{equation}
C(d)=\frac{\|v_{1:d}\|^2}{\|v\|^2}
\end{equation}
where $v_{1:d}$ is the prefix of $v$.

\subsection{Variance Concentration and Anisotropy (RQ3)}
\begin{equation}
C_{var}(d)=\frac{\mathrm{Var}(\langle f^{(d)}(x),v_{1:d}\rangle)}{\mathrm{Var}(\langle f^{(D)}(x),v\rangle)}
\end{equation}
with bootstrap confidence intervals.

\begin{equation}
\alpha(d)=\frac{C_{var}(d)}{\|v_{1:d}\|^2}
\end{equation}
Interpretation: $\alpha=1$ isotropic expression, $\alpha<1$ underexpression, $\alpha>1$ overexpression.
Because $\|v\|=1$, the term $\|v_{1:d}\|^2$ is the geometric share of bias-direction energy in the prefix. Therefore, $\alpha(d)=1$ means observed variance concentration exactly matches this geometric expectation (isotropic expression), while $\alpha(d)<1$ indicates suppressed variance expression relative to geometry.

\section{Experimental Setup}
\textbf{Models.} MPNet-Standard, MPNet-MRL, Nomic-MRL, mxbai-MRL, Qwen3-0.6B-MRL, Qwen3-4B-MRL.\\
\textbf{Dimensions.} Model-specific Matryoshka prefixes (e.g., 64/128/256/512/768 for 768-d models; 64/128/256/512/1024 for 1024-d models; plus 2560 for Qwen3-4B).\\
\textbf{Uncertainty.} Bootstrap CIs for $C_{var}(d)$ and induced CIs for $\alpha(d)$.\\
\textbf{Reproducibility.} Canonical evidence is frozen to generated JSON/plot artifacts listed in Section~\ref{sec:evidence}.

\section{Results}
\subsection{Layer 1: MPNet baseline (RQ1--RQ2)}
MPNet-Standard and MPNet-MRL show dimension-dependent shifts in WEAT/SEAT/StereoSet/CrowS and geometric concentration $C(d)$.

\subsection{Layer 2: Cross-model generalization}
Nomic-MRL, mxbai-MRL, and Qwen variants validate that dimension-wise bias behavior is model-dependent but nontrivial across architectures/scales.

\subsection{Layer 3: Main contribution (RQ3)}
$C_{var}(d)$ and $\alpha(d)$ reveal that several MRL models underexpress early-dimension bias variance despite geometric concentration evidence.

\begin{table}[t]
\caption{Anisotropy summary from frozen artifact \texttt{step13d\_alpha\_results.json} (gender axis).}
\label{tab:alpha-summary}
\centering
\begin{tabular}{lcc}
\toprule
Model & $\min\,\alpha(d)$ & Dimension \\
\midrule
MPNet-Standard & 0.727 & 256 \\
MPNet-MRL & 0.429 & 128 \\
Nomic-MRL & 0.553 & 128 \\
mxbai-MRL & 0.820 & 512 \\
Qwen3-0.6B-MRL & 0.358 & 64 \\
Qwen3-4B-MRL & 0.774 & 256 \\
\bottomrule
\end{tabular}
\end{table}

\section{Discussion}
The key measurement insight is that geometric and statistical views diverge under MRL truncation. Therefore, fairness claims based only on geometric concentration or only on one benchmark score are incomplete. For deployment, dimension-adaptive systems should consider both concentration and expression when selecting truncation points.

\section{Limitations and Ethics}
This version is gender-primary and English-centric. Race findings are exploratory unless the full $C_{var}/\alpha$ pipeline is replicated for race (and ideally additional social axes). Benchmark artifacts and lexical choices may encode historical harms; results should be interpreted as measurement outcomes, not normative statements.

\section{Conclusion}
MRL changes not only \emph{where} bias energy sits in vector space but also \emph{how strongly} that bias is statistically expressed across dimensions. The anisotropy factor $\alpha(d)$ provides a practical bridge between geometric intuition and variance-based fairness analysis.

\section{Canonical Evidence Map}
\label{sec:evidence}
Paths below are literal repository paths (legacy folder names preserved as-is).
\begin{table*}[t]
\caption{Frozen evidence artifacts used for paper claims.}
\centering
\begin{tabular}{lll}
\toprule
Layer & Purpose & Artifact paths (repository-relative) \\
\midrule
Layer 1 & MPNet baseline metrics + $C(d)$ & \texttt{step3\_weat.py}, \texttt{step4\_subspace.py}, \texttt{step5\_seat.py}, \texttt{step5\_stereoset\_v2.py}, \texttt{step6\_crowspairs\_new.py}, \texttt{step7\_all\_results.json} \\
Layer 2 & Cross-model generalization & \texttt{nomic/step8\_nomic\_results.json}, \texttt{mixbai\_mrl/step9\_mxbai\_results.json} (mxbai model), \texttt{qwen\_0.6\_emb.py/step10\_qwen3\_results.json} (Qwen3-0.6B folder), \texttt{qwen-3b/step11\_qwen3\_4b\_results.json} \\
Layer 3 & Variance + anisotropy contribution & \texttt{cvar\_new/step13b\_cvar\_results\_with\_ci.json}, \texttt{step13d\_alpha\_results.json}, \texttt{cvar\_new/Publication\_Quality\_C\_var(d)\_Plot\_with\_Bootstrap\_CIs.py}, \texttt{anistropy\_plots.py} \\
\bottomrule
\end{tabular}
\end{table*}

\section*{Figure and Table Placement Plan}
\begin{itemize}
    \item \textbf{Fig. 1 (main):} MPNet Standard vs MPNet-MRL prefix bias overview.
    \item \textbf{Fig. 2 (main):} Cross-model $C(d)$ comparison.
    \item \textbf{Fig. 3 (main):} $C_{var}(d)$ with bootstrap CIs.
    \item \textbf{Fig. 4 (main):} $\alpha(d)$ curves (headline figure).
    \item \textbf{Tab. 1 (main):} Model and dimension protocol.
    \item \textbf{Tab. 2 (main):} Compact metric summary across models.
    \item \textbf{Tab. 3 (main):} $\alpha$ minima and dimensions (Table~\ref{tab:alpha-summary}).
\end{itemize}

\section*{Appendix Plan}
\begin{itemize}
    \item Additional WEAT variants and full per-dimension score tables.
    \item StereoSet/CrowS detailed tables.
    \item Exploratory race analysis from early-layer metrics.
    \item Additional generated plots not selected for the main narrative.
\end{itemize}

\begin{thebibliography}{00}
\bibitem{caliskan2017} A.~Caliskan, J.~J. Bryson, and A.~Narayanan, ``Semantics derived automatically from language corpora contain human-like biases,'' \emph{Science}, 2017.
\bibitem{may2019} C.~May, A.~Wang, S.~Bordia, S.~Bowman, and R.~Rudinger, ``On measuring social biases in sentence encoders,'' \emph{NAACL}, 2019.
\bibitem{nangia2020} N.~Nangia, C.~Vania, R.~Bhalerao, and S.~R. Bowman, ``CrowS-Pairs: A challenge dataset for measuring social biases in masked language models,'' \emph{EMNLP}, 2020.
\bibitem{nadeem2021} M.~Nadeem, A.~Bethke, and S.~Reddy, ``StereoSet: Measuring stereotypical bias in pretrained language models,'' \emph{ACL}, 2021.
\bibitem{kusupati2022} A.~Kusupati et al., ``Matryoshka representation learning,'' \emph{NeurIPS}, 2022.
\end{thebibliography}

\end{document}
